% ============================================================
% OP_GUT1_recon_v1.tex  (PROPOSAL-ONLY RECONNAISSANCE)
% OP-GUT1 / SU(3): map of existing material, overclaim risks,
% candidate decomposition, formulation variants.
% NO implementation, NO derivational claims. Per GPT mandate.
% ============================================================
\section*{OP-GUT1 reconnaissance (proposal-only)}

\subsection*{0. Map of existing material (statuses as in the text)}
\begin{itemize}
  \item \textbf{Obstruction A} (triplet-condensate route):
    condensate in $\mathbf 3$ of $SU(3)$ has stabiliser $SU(2)$;
    exact colour cannot be carried by the condensate.
    Proof-sketch present; reads as a Level~A structural no-go.
  \item \textbf{Obstruction B} (dimensional exclusion):
    $\dim\mathfrak{su}(3)=8>\dim\mathfrak{so}(4)=6$, so
    $\mathfrak{su}(3)\not\subset\mathfrak{so}(4)$; maximal gauge
    structure from $O(4)\to O(3)$ is
    $U(2)=(SU(2)\times U(1))/\mathbb Z_2$. Level~A algebraic fact.
  \item \textbf{P7} (reverse-engineered postulate): local complex
    rank-3 module $\mathcal C_x\simeq\mathbb C^3$ with Hermitian
    form $h$ and complex volume form
    $\Omega\in\Lambda^3\mathcal C_x^*$; condensate singlet; frame
    redundancy $SU(3)$. Explicitly \emph{not} derived.
  \item \textbf{$U(3)\to SU(3)$ reduction} via $\Omega$
    ($\varepsilon_{abc}$): part of the candidate completion.
  \item \textbf{Zero-mode closure programme}
    (eq.~p7\_zero\_mode\_target; closure-targets paragraph):
    $\mathcal C_x\simeq\ker\mathcal D_{\rm defect}$,
    $\dim_{\mathbb C}=3$, $(h,\Omega)$ induced, colour connection
    as Berry--Wilczek--Zee connection. Recorded target; Open.
  \item \textbf{$G_2$ candidate} (nested, deeper, not adopted):
    $\mathrm{Stab}_{G_2}(u)=SU(3)$ (eq.~G2\_stab),
    $J_u(v)=u\times v$, $u^\perp\simeq\mathbb C^3$ (eq.~G2\_J).
    Standard mathematics, currently \emph{uncited}.
  \item \textbf{Minimal colour EFT under P7}: Level~B; admissible
    operator class; background-level triviality (clean negative).
  \item \textbf{Registry}: OP-GUT1 status ``Reframed'' (under P7
    with the zero-mode target); $N_{\rm colour}=3$ = minimal
    closure hypothesis (Level~C); OP-GUT2 (fermion reps) Open;
    anomaly side now carries the AN1--AN3 package, including the
    $SU(3)^3$ vector-like forcing remark.
\end{itemize}

\subsection*{1. What ``deriving $SU(3)$'' actually decomposes into}
(i)~existence of a rank-3 complex internal module from condensate
dynamics (the zero-mode theorem) --- the hard dynamical question;
(ii)~origin of the volume form $\Omega$, i.e.\ why $SU(3)$ rather
than $U(3)$;
(iii)~rank selection: why $N_c=3$;
(iv)~origin of Yang--Mills kinetics for the connection (partially
addressed by the EFT section; dynamical origin open).
These are logically independent; conflating them is the main
overclaim hazard of this OP.

\subsection*{2. Candidate decomposition (no commitments)}
\textbf{GUT1.0 --- status-audit and citation pass (low risk,
tractable now).}
(a)~Add literature citations for the $G_2$ facts
(candidates, to be page-verified before insertion:
Harvey--Lawson, \emph{Calibrated geometries}, Acta Math.\ 1982;
Baez, \emph{The octonions}, Bull.\ AMS 2002; possibly
Bryant 1987); the facts are standard, the article currently states
them bare.
(b)~Ensure the two obstructions carry explicit Level~A markers in
the local status table (echo pass only).

\textbf{GUT1.1 --- $N_c$-selection classification (candidate
Level~A arithmetic under inherited assumptions; the AN3 analogue).}
Generalise the anomaly system to $N_c$ colours with the inherited
electroweak content and an inherited charge-integrality /
atom-neutrality input; classify which $N_c$ are admissible.
Known literature direction (to be verified page-level before any
use: Abbas~1990-type analyses of charge quantisation with arbitrary
$N_c$).
Honest outcome shape: ``under the stated inherited inputs the
anomaly system pins $N_c=3$'' --- a conditional classification that
would sharpen the $N_{\rm colour}=3$ row from hypothesis~(C) to a
conditional result, \emph{not} a derivation of colour.
Moderate risk; requires own algebra + verified references.

\textbf{GUT1.2 --- $G_2$ mathematical kernel (candidate Level~A
conditional mathematics; presentational risk).}
State $\mathrm{Stab}_{G_2}(u)=SU(3)$ and the induced
$(J_u,h,\Omega)$ on $u^\perp$ as cited mathematical facts with a
short proof sketch, upgrading the chain
``$G_2$ structure $\Rightarrow$ the $SU(3)$-structure data of P7''
to Level~A \emph{conditional} mathematics.
The hypothesis (internal octonionic fibre) remains open/speculative
and the primary completion remains primary --- wording must keep
the conditional firewall.

\textbf{GUT1.3 --- zero-mode programme formalisation (Open;
definition-level).}
Sharpen eq.~p7\_zero\_mode\_target into a precise problem
statement: axioms on $\mathcal D_{\rm defect}$ (locality,
condensate equivariance, ellipticity class), the exact meaning of
``$(h,\Omega)$ induced'', and the candidate route taxonomy
(index/defect-topology count; three-fold internal degeneracy; $G_2$
reduction).
No toy model, no existence claim; output is a definition + open
problem block.

\textbf{GUT1.4 --- untouched open remainder.}
Yang--Mills kinetic origin, $g_s$, confinement, OP-GUT2
representations, $\theta_{\rm QCD}$. No new claims.

\subsection*{3. Overclaim risk register}
Forbidden:
``ECT derives $SU(3)$'';
``$N_c=3$ is derived'';
``the medium carries a $G_2$ structure'';
``P7 is now derived'';
inverting the obstruction theorems into positive claims;
presenting the zero-mode target as partially achieved.
Safe:
``structure-preserving frame redundancy under P7'';
``conditional mathematical chain ($G_2\Rightarrow$ P7 data)'';
``$N_c$ pinned under stated inherited inputs'';
``programme specification, not existence''.

\subsection*{4. Formulation variants (for the upgradable bits only)}
(a)~$G_2$ facts: ``standard mathematical facts
[cites]; within ECT they function only as the conditional kernel of
the non-adopted deeper candidate.''
(b)~$N_c$ row (if GUT1.1 succeeds):
``$N_{\rm colour}=3$ --- conditional classification: pinned by the
anomaly system under inherited electroweak content and charge
integrality (\S\dots); origin of the colour module itself remains
open (P7/zero-mode programme).''
(c)~Zero-mode block heading: ``Definition (colour zero-mode
problem)'' + ``Open problem (OP-GUT1, sharpened)''.

\subsection*{5. Recommended sequencing}
GUT1.0 (citations + status echoes) $\to$ GUT1.1 (only after
page-level literature verification) $\to$ GUT1.3 (formalisation)
$\to$ GUT1.2 (optional, only with the conditional firewall).
GUT1.1 is the single genuinely new mathematical content with
AN3-like safety properties; everything else is consolidation.

\subsection*{6. Questions for GPT}
(Q1)~Approve the decomposition GUT1.0--GUT1.4?
(Q2)~Does the $N_c$-from-anomalies classification belong under
OP-GUT1 or under the anomaly family (OP-GUT6 vicinity)?
(Q3)~$G_2$ citations: agree to add, and to keep GUT1.2 strictly
conditional (or drop GUT1.2 entirely)?
(Q4)~Zero-mode formalisation: place as a definition-level block
inside \S colour\_completion, or as an expansion of the
closure-targets paragraph?
(Q5)~Any objection to the Abbas-type literature route for GUT1.1,
pending page-level verification?
